\subsection{Parametrické vyjádření roviny}
\label{subsec:Parametric_plane}

Rovina je určena jedním bodem a dvěma navzájem nezávislými směry.
Z výchozího bodu se tedy můžeme pohybovat současně ve směru dvou
nenulových, navzájem nerovnoběžných vektorů.

Uvažujme bod
\[
    A[x_A;y_A;z_A]
\]
a dva nenulové, navzájem nerovnoběžné vektory
\[
    \vec{u}=(u_x;u_y;u_z),
    \qquad
    \vec{v}=(v_x;v_y;v_z).
\]

Libovolný bod $X[x;y;z]$ roviny získáme z bodu $A$ posunutím o vhodný
násobek vektoru $\vec{u}$ a vhodný násobek vektoru $\vec{v}$:
\[
    X=A+s\vec{u}+t\vec{v},
\]
kde $s,t\in\mathbb{R}$.


\subsubsection{Rovina určená bodem a dvěma směrovými vektory}
\label{subsubsec:Plane_point_two_direction_vectors}

\begin{definition}
    \label{def:Parametric_plane}

    Nechť $A[x_A;y_A;z_A]$ je bod v prostoru a nechť
    \[
        \vec{u}=(u_x;u_y;u_z),
        \qquad
        \vec{v}=(v_x;v_y;v_z)
    \]
    jsou dva nenulové, navzájem nerovnoběžné vektory.

    Rovina $\alpha$ procházející bodem $A$ a mající směrové vektory
    $\vec{u}$ a $\vec{v}$ je množina všech bodů $X$, které lze vyjádřit
    ve tvaru
    \begin{equation}
        X=A+s\vec{u}+t\vec{v},
        \qquad
        s,t\in\mathbb{R}.
        \label{eq:Parametric_plane_vector}
    \end{equation}

    Po rozepsání do souřadnic získáme parametrické rovnice roviny
    \begin{equation}
        \begin{aligned}
            \alpha: x&=x_A+s u_x+t v_x,\\
            y&=y_A+s u_y+t v_y,\\
            z&=z_A+s u_z+t v_z; \quad s,t\in\mathbb{R}.
        \end{aligned}
        \label{eq:Parametric_plane_coordinates}
    \end{equation}

    Parametry $s$ a $t$ jsou opět známými \textbf{parametry roviny} a vektory
    $\vec{u}$ a $\vec{v}$ jejími \textbf{směrovými vektory}.
\end{definition}

Vektor směřující z bodu $A$ do libovolného bodu $X$ roviny má tvar
\begin{equation}
    \overrightarrow{AX}=s\vec{u}+t\vec{v}.
    \label{eq:Plane_displacement_vector}
\end{equation}

Každý vektor spojující bod $A$ s bodem roviny je tedy lineární kombinací
jejích směrových vektorů.

\begin{figure}[H]
    \centering
    \begin{tikzpicture}[scale=0.9]

        % Rovina
        \coordinate (P1) at (-1.2,0.1,0.3);
        \coordinate (P2) at (4.4,1.2,0.9);
        \coordinate (P3) at (5.0,3.8,2.9);
        \coordinate (P4) at (-0.6,2.7,2.3);

        \fill[gray!18] (P1) -- (P2) -- (P3) -- (P4) -- cycle;
        \draw[thin] (P1) -- (P2) -- (P3) -- (P4) -- cycle;

        % Osy - kladná část
        \draw[->] (0,0,0) -- (5.2,0,0) node[right] {$x$};
        \draw[->] (0,0,0) -- (0,4.7,0) node[above] {$y$};
        \draw[->] (0,0,0) -- (0,0,4.1) node[below left] {$z$};

        % Osy - záporná část
        \draw[->, dashed] (0,0,0) -- (-2.2,0,0);
        \draw[->, dashed] (0,0,0) -- (0,-2.2,0);
        \draw[->, dashed] (0,0,0) -- (0,0,-2.7);

        \node[above left] at (0,0,0) {$O$};

        % Body
        \coordinate (A) at (1,1,1);
        \coordinate (U) at (3,1.4,1.2);
        \coordinate (V) at (1.6,2.4,2);
        \coordinate (X) at (3.6,2.8,2.2);

        % Směrové vektory
        \draw[vector] (A) -- (U)
            node[midway, sloped, above] {$\vec{u}$};

        \draw[vector] (A) -- (V)
            node[midway, sloped, above] {$\vec{v}$};

        % Pomocný rovnoběžník
        \draw[dashed] (U) -- (X);
        \draw[dashed] (V) -- (X);

        % Body
        \fill (A) circle (2pt) node[below left] {$A$};
        \fill (X) circle (2pt) node[above right] {$X$};

        \node at (4.35,3,2.65) {$\alpha$};

    \end{tikzpicture}

    \caption{Bod $X$ roviny $\alpha$ získáme z bodu $A$ posunutím
    ve směrech vektorů $\vec{u}$ a $\vec{v}$.}
    \label{fig:Parametric_plane_geometric_meaning}
\end{figure}

Na obrázku~\ref{fig:Parametric_plane_geometric_meaning} je znázorněn
geometrický význam vztahu~\eqref{eq:Parametric_plane_vector}.
Posunutí
\[
    s\vec{u}+t\vec{v}
\]
může mít libovolnou velikost i orientaci v rámci roviny určené směry
$\vec{u}$ a $\vec{v}$.


\subsubsection{Význam parametrů}
\label{subsubsec:Meaning_of_plane_parameters}

Parametry $s$ a $t$ se mohou měnit nezávisle na sobě. Každá uspořádaná
dvojice
\[
    (s;t)\in\mathbb{R}^2
\]
určuje právě jeden bod roviny.

Pro některé zvláštní hodnoty parametrů dostáváme:

\begin{itemize}
    \item pro $s=0$ a $t=0$ získáme výchozí bod
    \[
        X=A;
    \]

    \item pro $t=0$ se pohybujeme pouze ve směru vektoru $\vec{u}$:
    \[
        X=A+s\vec{u},
    \]
    dostáváme tím tedy přímku procházející bodem $A$ a rovnoběžnou s vektorem $\vec{u}$;

    \item pro $s=0$ se pohybujeme pouze ve směru vektoru $\vec{v}$:
    \[
        X=A+t\vec{v},
    \]
    dostáváme tím tedy přímku procházející bodem $A$ a rovnoběžnou s vektorem $\vec{v}$;

    \item pro pevně zvolené $t$ a proměnné $s$ získáváme přímku
    rovnoběžnou s vektorem $\vec{u}$;

    \item pro pevně zvolené $s$ a proměnné $t$ získáváme přímku
    rovnoběžnou s vektorem $\vec{v}$.
\end{itemize}

Rovinu si tak můžeme představit jako nekonečnou soustavu navzájem
rovnoběžných přímek. Například při pevně zvoleném parametru $t$
a proměnném parametru $s$ získáváme přímky rovnoběžné se směrovým
vektorem $\vec{u}$, jejichž sjednocením je celá rovina.

\begin{remark}
    Pro parametrické vyjádření přímky stačil jeden parametr, protože
    přímka má pouze jeden nezávislý směr. Rovina má dva nezávislé
    směry, a proto potřebujeme dva parametry.
\end{remark}


\subsubsection{Proč musí být směrové vektory nerovnoběžné}
\label{subsubsec:Why_plane_vectors_nonparallel}

Podmínka, že vektory $\vec{u}$ a $\vec{v}$ nejsou rovnoběžné, je
nezbytná.

Kdyby například existovalo číslo $k\in\mathbb{R}$ takové, že
\[
    \vec{v}=k\vec{u},
\]
potom by platilo
\begin{align*}
    X
    &=A+s\vec{u}+t\vec{v}\\
    &=A+s\vec{u}+tk\vec{u}\\
    &=A+(s+kt)\vec{u}.
\end{align*}

Místo dvou nezávislých směrů bychom tedy získali pouze násobky jediného
vektoru $\vec{u}$. Výslednou množinou bodů by nebyla rovina, ale přímka.

Podobný problém nastane, pokud je některý ze směrových vektorů nulový.
Nulový vektor neurčuje žádný směr a jeho násobky nemění polohu bodu.

\begin{remark}
    Dva vektory mohou tvořit dvojici směrových vektorů roviny právě
    tehdy, jsou-li lineárně nezávislé. Pro dvojici vektorů to znamená,
    že ani jeden z nich není reálným násobkem druhého.
\end{remark}


\subsubsection{Sestavení parametrického vyjádření roviny}
\label{subsubsec:Constructing_parametric_plane}

Při sestavování parametrického vyjádření roviny potřebujeme:

\begin{enumerate}
    \item jeden bod $A$, který v rovině leží;
    \item dva nenulové, navzájem nerovnoběžné směrové vektory
    $\vec{u}$ a $\vec{v}$ ležící v rovině;
    \item dosadit tyto údaje do vztahu
    \[
        X=A+s\vec{u}+t\vec{v}.
    \]
\end{enumerate}

\begin{example}
    \label{exa:Parametric_plane_point_vectors}

    Určete parametrické vyjádření roviny $\alpha$, která prochází bodem
    $A[1;-2;3]$ a má směrové vektory $\vec{u}=(2;1;-1)\text{, } \vec{v}=(1;-1;2)$.

    Vektory $\vec{u}$ a $\vec{v}$ nejsou rovnoběžné, a proto určují
    dva nezávislé směry roviny.

    Použijeme vztah~\eqref{eq:Parametric_plane_vector}:
    \begin{align*}
        X
        &=A+s\vec{u}+t\vec{v}\\
        &=[1;-2;3]+s(2;1;-1)+t(1;-1;2)\\
        &=[1+2s+t;-2+s-t;3-s+2t].
    \end{align*}

    Parametrické vyjádření roviny je
    \[
        \begin{aligned}
            \alpha: x&=1+2s+t,\\
            y&=-2+s-t,\\
            z&=3-s+2t; \quad s,t\in\mathbb{R}.
        \end{aligned}
    \]
\end{example}

Z parametrického vyjádření můžeme vytvářet jednotlivé body roviny
volbou konkrétních hodnot parametrů. Pro rovinu z předchozího příkladu
například dostáváme

\[
    \begin{array}{c|c|c}
        s & t & X[x;y;z]\\
        \hline
        0 & 0 & A[1;-2;3]\\
        1 & 0 & B[3;-1;2]\\
        0 & 1 & C[2;-3;5]\\
        1 & 1 & D[4;-2;4]
    \end{array}
\]

Body $A$, $B$, $C$ a $D$ tvoří rovnoběžník. Platí totiž
\[
    \overrightarrow{AB}=\vec{u},
    \qquad
    \overrightarrow{AC}=\vec{v},
\]
a bod $D$ získáme současným posunutím o oba směrové vektory:
\[
    D=A+\vec{u}+\vec{v}.
\]


\subsubsection{Čtení parametrického vyjádření roviny}
\label{subsubsec:Reading_parametric_plane}

Z parametrických rovnic můžeme přímo vyčíst jeden bod roviny a dva
její směrové vektory. Konstantní členy určují bod odpovídající hodnotám
$s=0$, $t=0$. Koeficienty u parametru $s$ tvoří souřadnice prvního
směrového vektoru a koeficienty u parametru $t$ souřadnice druhého.

\begin{example}
    \label{exa:Reading_parametric_plane}

    Je dána rovina
    \[
        \begin{aligned}
            \beta: x&=-1+3s-t,\\
            y&=2-s+2t,\\
            z&=4+2s+t; \quad s,t\in\mathbb{R}.
        \end{aligned}
    \]

    Určete jeden bod roviny a dva její směrové vektory.

    Pro $s=0$ a $t=0$ dostáváme bod
    \[
        K[-1;2;4].
    \]

    Koeficienty u parametru $s$ tvoří směrový vektor
    \[
        \vec{u}=(3;-1;2),
    \]
    zatímco koeficienty u parametru $t$ tvoří druhý směrový vektor
    \[
        \vec{v}=(-1;2;1).
    \]

    Rovina $\beta$ tedy prochází bodem $K[-1;2;4]$ a má směrové vektory
    \[
        \vec{u}=(3;-1;2),
        \qquad
        \vec{v}=(-1;2;1).
    \]
\end{example}

Při čtení směrových vektorů je třeba dávat pozor zejména na znaménka.
Například z rovnice
\[
    y=2-s+2t
\]
čteme jako druhou souřadnici prvního směrového vektoru číslo $-1$
a jako druhou souřadnici druhého směrového vektoru číslo $2$.


\subsubsection{Příslušnost bodu k rovině}
\label{subsubsec:Point_membership_parametric_plane}

Aby bod $P[x_P;y_P;z_P]$ ležel v parametricky vyjádřené rovině,
musí existovat jedna dvojice reálných parametrů $(s;t)$, která
současně vytvoří všechny tři souřadnice bodu $P$.

Pro rovinu
\[
    \begin{aligned}
        \alpha: x&=x_A+s u_x+t v_x,\\
        y&=y_A+s u_y+t v_y,\\
        z&=z_A+s u_z+t v_z; \quad s,t\in\mathbb{R}
    \end{aligned}
\]
tedy platí
\[
    P\in\alpha
    \quad\Longleftrightarrow\quad
    \text{existují }s,t\in\mathbb{R}\text{ takové, že}
    \begin{cases}
        x_P=x_A+s u_x+t v_x,\\
        y_P=y_A+s u_y+t v_y,\\
        z_P=z_A+s u_z+t v_z.
    \end{cases}
\]

Nestačí, aby některé dvě rovnice určily dvojici parametrů. Stejná
dvojice hodnot $s$ a $t$ musí splňovat všechny tři rovnice.

\begin{example}
    \label{exa:Point_membership_parametric_plane}

    Rozhodněte, zda body
    \[
        P[4;-2;4],
        \qquad
        Q[3;-1;3]
    \]
    leží v rovině
    \[
        \begin{aligned}
            \alpha: x&=1+2s+t,\\
            y&=-2+s-t,\\
            z&=3-s+2t; \quad s,t\in\mathbb{R}.
        \end{aligned}
    \]

    Nejprve dosadíme souřadnice bodu $P[4;-2;4]$:
    \[
        \begin{cases}
            4=1+2s+t,\\
            -2=-2+s-t,\\
            4=3-s+2t.
        \end{cases}
    \]

    Po úpravě
    \[
        \begin{cases}
            2s+t=3,\\
            s-t=0,\\
            -s+2t=1.
        \end{cases}
    \]

    Z druhé rovnice plyne
    \[
        s=t.
    \]

    Po dosazení do první rovnice získáme
    \[
        3s=3,
    \]
    tedy
    \[
        s=1,
        \qquad
        t=1.
    \]

    Ověříme třetí rovnici:
    \[
        -1+2\cdot1=1.
    \]

    Stejná dvojice parametrů $(s;t)=(1;1)$ splňuje všechny tři
    rovnice, a proto
    \[
        P\in\alpha.
    \]

    Nyní dosadíme souřadnice bodu $Q[3;-1;3]$:
    \[
        \begin{cases}
            3=1+2s+t,\\
            -1=-2+s-t,\\
            3=3-s+2t.
        \end{cases}
    \]

    Po úpravě
    \[
        \begin{cases}
            2s+t=2,\\
            s-t=1,\\
            -s+2t=0.
        \end{cases}
    \]

    Z prvních dvou rovnic dostáváme
    \[
        s=1,
        \qquad
        t=0.
    \]

    Tato dvojice však nesplňuje třetí rovnici, protože
    \[
        -1+2\cdot0=-1\neq0.
    \]

    Proto
    \[
        Q\notin\alpha.
    \]
\end{example}

Příslušnost bodu můžeme vyjádřit také vektorově. Bod $P$ leží
v rovině procházející bodem $A$ se směrovými vektory $\vec{u}$
a $\vec{v}$ právě tehdy, když lze vektor
\[
    \overrightarrow{AP}=P-A
\]
vyjádřit jako lineární kombinaci
\[
    \overrightarrow{AP}=s\vec{u}+t\vec{v}
\]
pro nějaká $s,t\in\mathbb{R}$.


\subsubsection{Rovina určená třemi body}
\label{subsubsec:Plane_through_three_points}

Rovina bývá často zadána třemi body
\[
    A[x_A;y_A;z_A],
    \qquad
    B[x_B;y_B;z_B],
    \qquad
    C[x_C;y_C;z_C].
\]

Aby tyto body určovaly právě jednu rovinu, nesmějí ležet na jedné
přímce.

Jako výchozí bod můžeme zvolit například bod $A$. Dva směrové vektory
roviny získáme jako
\[
    \overrightarrow{AB}
    =
    (x_B-x_A;y_B-y_A;z_B-z_A)
\]
a
\[
    \overrightarrow{AC}
    =
    (x_C-x_A;y_C-y_A;z_C-z_A).
\]

Pokud body $A$, $B$ a $C$ neleží na jedné přímce, nejsou vektory
$\overrightarrow{AB}$ a $\overrightarrow{AC}$ rovnoběžné a mohou
sloužit jako směrové vektory roviny.

Parametrické vyjádření potom má tvar
\begin{equation}
    X=A+s\overrightarrow{AB}+t\overrightarrow{AC},
    \qquad
    s,t\in\mathbb{R}.
    \label{eq:Plane_through_three_points}
\end{equation}

Po rozepsání do souřadnic:
\begin{equation}
    \begin{aligned}
        \alpha: x&=x_A+s(x_B-x_A)+t(x_C-x_A),\\
        y&=y_A+s(y_B-y_A)+t(y_C-y_A),\\
        z&=z_A+s(z_B-z_A)+t(z_C-z_A);
        \quad s,t\in\mathbb{R}.
    \end{aligned}
    \label{eq:Plane_through_three_points_coordinates}
\end{equation}

\begin{example}
    \label{exa:Plane_through_three_points}

    Určete parametrické vyjádření roviny $\alpha$, která prochází body
    \[
        A[1;-2;3],
        \qquad
        B[3;-1;2],
        \qquad
        C[2;-3;5].
    \]

    Jako výchozí bod zvolíme bod $A$. Určíme vektory
    \[
        \overrightarrow{AB}
        =
        (3-1;-1-(-2);2-3)
        =
        (2;1;-1)
    \]
    a
    \[
        \overrightarrow{AC}
        =
        (2-1;-3-(-2);5-3)
        =
        (1;-1;2).
    \]

    Vektory $\overrightarrow{AB}$ a $\overrightarrow{AC}$ nejsou
    rovnoběžné, takže zadané body neleží na jedné přímce a určují
    právě jednu rovinu.

    Parametrické vyjádření roviny je
    \[
        \begin{aligned}
            \alpha: x&=1+2s+t,\\
            y&=-2+s-t,\\
            z&=3-s+2t; \quad s,t\in\mathbb{R}.
        \end{aligned}
    \]
\end{example}

\begin{remark}
    Jako výchozí bod nemusíme zvolit právě bod $A$. Můžeme použít
    také bod $B$ nebo $C$ a vytvořit dva vhodné vektory vycházející
    z tohoto bodu. Výsledné parametrické vyjádření bude vypadat jinak,
    ale bude popisovat tutéž rovinu.
\end{remark}

Pokud by byly vektory $\overrightarrow{AB}$ a $\overrightarrow{AC}$
rovnoběžné, ležely by body $A$, $B$ a $C$ na jedné přímce. Takovými
třemi body prochází nekonečně mnoho rovin, a proto by rovina nebyla
určena jednoznačně.


\subsubsection{Rovina určená přímkou a bodem}
\label{subsubsec:Plane_line_and_point}

Rovinu lze jednoznačně určit také přímkou $p$ a bodem $C$, který
na této přímce neleží.

Nechť je přímka zadána parametricky ve tvaru
\[
    X=A+r\vec{u},
    \qquad
    r\in\mathbb{R}.
\]

Jeden směrový vektor roviny získáme přímo z přímky:
\[
    \vec{u}.
\]

Druhý směrový vektor vytvoříme spojením bodu $A$ přímky s bodem $C$:
\[
    \overrightarrow{AC}=C-A.
\]

Protože bod $C$ na přímce neleží, nejsou vektory $\vec{u}$ a
$\overrightarrow{AC}$ rovnoběžné.

Parametrické vyjádření roviny má tedy tvar
\begin{equation}
    X=A+s\vec{u}+t\overrightarrow{AC},
    \qquad
    s,t\in\mathbb{R}.
    \label{eq:Plane_line_and_point}
\end{equation}

\begin{example}
    \label{exa:Plane_line_and_point}

    Určete parametrické vyjádření roviny $\alpha$, která obsahuje přímku
    \[
        \begin{aligned}
            p: x&=1+2r,\\
            y&=-2+r,\\
            z&=3-r; \quad r\in\mathbb{R}
        \end{aligned}
    \]
    a bod
    \[
        C[2;-3;5].
    \]

    Z parametrického vyjádření přímky vyčteme bod
    \[
        A[1;-2;3]
    \]
    a její směrový vektor
    \[
        \vec{u}=(2;1;-1).
    \]

    Druhý směrový vektor roviny získáme jako
    \[
        \overrightarrow{AC}
        =
        (2-1;-3-(-2);5-3)
        =
        (1;-1;2).
    \]

    Parametrické vyjádření roviny je
    \[
        \begin{aligned}
            \alpha: x&=1+2s+t,\\
            y&=-2+s-t,\\
            z&=3-s+2t; \quad s,t\in\mathbb{R}.
        \end{aligned}
    \]
\end{example}

\begin{remark}
    Kdyby bod $C$ ležel na přímce $p$, byl by vektor
    $\overrightarrow{AC}$ rovnoběžný se směrovým vektorem přímky.
    Získali bychom pouze jeden nezávislý směr a rovina by nebyla
    určena jednoznačně. Danou přímkou totiž prochází nekonečně mnoho
    rovin.
\end{remark}


\subsubsection{Rovina určená dvěma přímkami}
\label{subsubsec:Plane_determined_by_two_lines}

Dvě přímky určují rovinu pouze v některých případech.

\paragraph{Dvě různoběžné přímky.}

Jestliže se přímky $p$ a $q$ protínají v bodě $A$ a mají navzájem
nerovnoběžné směrové vektory $\vec{u}$ a $\vec{v}$, určují právě
jednu rovinu:
\[
    X=A+s\vec{u}+t\vec{v},
    \qquad
    s,t\in\mathbb{R}.
\]

\begin{example}
    \label{exa:Plane_two_intersecting_lines}

    Určete parametrické vyjádření roviny $\alpha$, která obsahuje
    přímky
    \[
        \begin{aligned}
            p: x&=1+r,\\
            y&=2-r,\\
            z&=3+2r; \quad r\in\mathbb{R}
        \end{aligned}
    \]
    a
    \[
        \begin{aligned}
            q: x&=1+2s,\\
            y&=2+s,\\
            z&=3-s; \quad s\in\mathbb{R}.
        \end{aligned}
    \]

    Pro $r=0$ a $s=0$ získáváme u obou přímek bod
    \[
        A[1;2;3].
    \]

    Přímky se tedy protínají v bodě $A$.

    Jejich směrové vektory jsou
    \[
        \vec{u}=(1;-1;2),
        \qquad
        \vec{v}=(2;1;-1).
    \]

    Vektory nejsou rovnoběžné, proto přímky určují právě jednu rovinu.

    Její parametrické vyjádření je
    \[
        X=A+k\vec{u}+l\vec{v},
    \]
    tedy
    \[
        \begin{aligned}
            \alpha: x&=1+k+2l,\\
            y&=2-k+l,\\
            z&=3+2k-l; \quad k,l\in\mathbb{R}.
        \end{aligned}
    \]
\end{example}

\paragraph{Dvě různé rovnoběžné přímky.}

Nechť přímky mají společný směrový vektor $\vec{u}$. Zvolíme bod $A$
na první přímce a bod $B$ na druhé přímce. Druhý směr roviny určí
vektor
\[
    \overrightarrow{AB}.
\]

Protože jsou přímky různé, není vektor $\overrightarrow{AB}$
rovnoběžný s vektorem $\vec{u}$. Rovina má parametrické vyjádření
\[
    X=A+s\vec{u}+t\overrightarrow{AB},
    \qquad
    s,t\in\mathbb{R}.
\]

\paragraph{Dvě mimoběžné přímky.}

Mimoběžné přímky neleží v jedné rovině. Žádná rovina je proto nemůže
současně obsahovat.

\begin{remark}
    Dvě přímky určují právě jednu rovinu tehdy, jsou-li různoběžné
    nebo jsou-li různé a rovnoběžné. Mimoběžné přímky společnou
    rovinu nemají a dvě totožné přímky leží v nekonečně mnoha rovinách.
\end{remark}


\subsubsection{Nejednoznačnost parametrického vyjádření roviny}
\label{subsubsec:Nonuniqueness_parametric_plane}

Stejně jako parametrické vyjádření přímky není ani parametrické
vyjádření roviny jednoznačné.

Pro tutéž rovinu můžeme zvolit:

\begin{itemize}
    \item libovolný její bod jako výchozí bod;
    \item libovolné dva lineárně nezávislé vektory ležící v této
    rovině jako směrové vektory;
    \item jiná označení parametrů.
\end{itemize}

Uvažujme rovinu
\[
    X=A+s\vec{u}+t\vec{v},
\]
kde
\[
    A[1;-2;3],
    \qquad
    \vec{u}=(2;1;-1),
    \qquad
    \vec{v}=(1;-1;2).
\]

Bod
\[
    B=A+\vec{u}=[3;-1;2]
\]
v této rovině leží.

Také vektor
\[
    \vec{w}
    =
    \vec{u}+\vec{v}
    =
    (3;0;1)
\]
je směrovým vektorem této roviny. Vektory $\vec{u}$ a $\vec{w}$
nejsou rovnoběžné, takže tutéž rovinu můžeme zapsat například jako
\[
    X=B+r\vec{u}+q\vec{w},
    \qquad
    r,q\in\mathbb{R}.
\]

Po rozepsání:
\[
    \begin{aligned}
        \alpha: x&=3+2r+3q,\\
        y&=-1+r,\\
        z&=2-r+q; \quad r,q\in\mathbb{R}.
    \end{aligned}
\]

Oba zápisy popisují tutéž rovinu, přestože používají jiný výchozí bod,
jinou dvojici směrových vektorů i jiná označení parametrů.

Skutečně platí
\begin{align*}
    B+r\vec{u}+q\vec{w}
    &=
    A+\vec{u}+r\vec{u}+q(\vec{u}+\vec{v})\\
    &=
    A+(1+r+q)\vec{u}+q\vec{v}.
\end{align*}

Každý bod druhého vyjádření je tedy bodem prvního vyjádření. Protože
hodnoty $1+r+q$ a $q$ mohou nabývat libovolných reálných hodnot, platí
také opačné tvrzení.

\begin{remark}
    Dva různé parametrické zápisy mohou popisovat tutéž rovinu.
    Nestačí proto porovnávat pouze výchozí body nebo jednotlivé
    směrové vektory. Rozhodující je, zda oba zápisy popisují stejnou
    množinu bodů a stejné dva nezávislé směry v prostoru.
\end{remark}


% ============================================================
%             P Ř E C H O D   K   O B E C N É   R O V N I C I
% ============================================================

Parametrické vyjádření roviny využívá dva navzájem nezávislé směry,
které v rovině leží. Rovinu však můžeme popsat také opačným způsobem:
pomocí směru, který je k celé rovině kolmý.

Má-li rovina směrové vektory
\[
    \vec{u}
    \qquad\text{a}\qquad
    \vec{v},
\]
potom jejich vektorový součin
\[
    \vec{n}
    =
    \vec{u}\times\vec{v}
\]
je kolmý současně k oběma směrovým vektorům, a tedy k celé rovině.

\begin{definition}
    Nenulový vektor $\vec{n}$ kolmý k rovině nazýváme
    \textbf{normálovým vektorem roviny}.
\end{definition}

Je-li
\[
    \vec{n}=(a;b;c),
\]
potom právě jeho souřadnice budou hrát zásadní roli při zápisu roviny
pomocí její \textbf{obecné rovnice}, kterou se budeme zabývat
v následující podkapitole.